% !TEX program = xelatex
\documentclass[9pt,aspectratio=169]{beamer}
\usetheme{metropolis}

% ── Packages ──────────────────────────────────────────────
\usepackage{amsmath,amssymb}
\usepackage{booktabs}
\usepackage{graphicx}
\usepackage{tikz}
\usepackage{appendixnumberbeamer}

% ── Figure paths ──────────────────────────────────────────
\graphicspath{%
  {../output/diagnostics/figures/}%
  {../output/did/figures/}%
  {../notes/empirics/compra_agil_descriptives/}%
}

% ── Table path helper ─────────────────────────────────────
% All \input{} paths are relative to this file's directory

% ── Meta ──────────────────────────────────────────────────
\title{Local Preference Policies in Public Procurement:\\
       Efficiency, Distributional Effects, and Spillovers}
\subtitle{Evidence from Chile's Compra Ágil Reform}
\author{Nicolás Jiménez}
\date{March 2026}
\institute{Yale University}

\begin{document}

% ══════════════════════════════════════════════════════════
% TITLE
% ══════════════════════════════════════════════════════════
\maketitle

% ══════════════════════════════════════════════════════════
% INTRODUCTION
% ══════════════════════════════════════════════════════════
\section{Introduction}

% ── Motivation ───────────────────────────────────────────
\begin{frame}{Motivation}
\textbf{[PLACEHOLDER — fill with motivation slide]}

\medskip
Suggested content:
\begin{itemize}
  \item Public procurement is large: OECD governments spend 12--15\% of GDP
        on procurement. Developing countries often more.
  \item Local preference policies are widespread: Buy American, EU local content
        rules, developing-country SME preferences.
  \item Tension: local preferences may protect domestic firms and create local
        jobs, but can reduce competition and raise costs for the government.
  \item Little causal evidence on the competition and welfare effects of these
        policies, especially in developing countries.
  \item Chilean Compra Ágil reform (Dec 2024) provides a sharp, quasi-random
        change in the scope of local preference.
\end{itemize}
\end{frame}

% ── Research Questions ───────────────────────────────────
\begin{frame}{Research Questions}
\begin{enumerate}
  \item[\textbf{Q1}] \textbf{Efficiency and distributional effects:} What are the short-run effects of local
        preference policies on competition in public procurement? 
        \begin{itemize}
          \item \textbf{Efficiency:} Number and composition of bidders; Government procurement costs (winning bid prices)
          \item[] $\to$ Do local entry gains offset non-local exclusion? 
          \item \textbf{Distributional effects:} Local vs.\ non-local firms; SME vs.\ large firms
        \end{itemize}

  \vspace{0.5em}
  \item[\textbf{Q2}] \textbf{Heterogeneity:} How do the effects of the policy vary by region, based on pre-existing market structure and exogenous characteristics?
        \begin{itemize}
          \item Heterogeneous by government demand, supply-side market thickness, distance from Santiago, and other characteristics
          \item Heterogeneity in efficiency and distributional implications of the policy
        \end{itemize}

  \vspace{0.5em}
\item[\textbf{Q3}] \textbf{Cross-market spillovers:} Do local preference policies for 30--100 UTM tenders cause them
      to redirect bids to untreated bands (100--500+ UTM), in their home
      regions or outside?
      \begin{itemize}
        \item \textbf{Mechanism:} firm-level capacity constraints linking
              treated and untreated markets across regions.
        \item \textbf{Outcomes in untreated bands:} entry, bidder composition, and winning prices.
        \item \textbf{Welfare:} do cross-market effects amplify or
              offset the direct efficiency and distributional effects in Q1?
      \end{itemize}
\end{enumerate}
\end{frame}



% ── Project Overview ─────────────────────────────────────
\begin{frame}{This Project}
\begin{columns}[T]
  \column{0.5\textwidth}
  \textbf{Part I: Empirics}
  \begin{itemize}
    \item Descriptive facts about local and non-local firm participation
          in Chilean public procurement
    \item DiD estimates of the short-run effects of the Compra Ágil reform
          on bidder entry, bidder composition, and winning prices
    \item Regional heterogeneity: how effects vary with pre-reform market
          structure and geography
    \item Cross-market spillovers: competition effects in \textit{above}-threshold
          markets driven by firm capacity reallocation
  \end{itemize}

  \column{0.5\textwidth}
  \textbf{Part II: Structural Model}
  \begin{itemize}
    \item Spatial model of firm entry and bidding for public procurement projects
    \item Firms have home regions; face distance-based cost and entry cost
          disadvantages; are subject to capacity constraints linking markets
    \item Estimation using pre- and post-reform variation
    \item \textbf{Counterfactuals:}
    \begin{itemize}
      \item[-] Optimal local preference threshold by region
      \item[-] Welfare decomposition (cost vs.\ local benefit)
      \item[-] Alternative policy designs
    \end{itemize}
  \end{itemize}
\end{columns}
\end{frame}

% ── Related Literature ───────────────────────────────────
\begin{frame}{Related Literature}
\textbf{[PLACEHOLDER — fill with related literature]}

\medskip
Key strands:
\begin{itemize}
  \item \textbf{Public procurement and competition:}
        Athey, Levin \& Seira (2011); Krasnokutskaya \& Seim (2011);
        Decarolis (2014); Coviello \& Mariniello (2014)
  \item \textbf{Local preference / Buy-national policies:}
        Marion (2007); Kang \& Miller (2022); Shingal (2020); Buy American paper;
  \item \textbf{SME preferences in procurement:}
        Carril, Gonzalez-Lira \& Walker (2022); Sorensen (2022)
  \item \textbf{Procurement reform in developing countries:}
        Bandiera et al.\ (2009); Best et al.\ (2017); Bosio et al.\ (2022)
  \item \textbf{Cross-market spillovers and capacity constraints:}
        Jeziorski \& Krasnokutskaya (2016); Hoekmann et al., (2025)
\end{itemize}

% \medskip
% \textbf{Contribution:} First quasi-experimental estimates of local preference
% effects accounting for cross-market spillovers via capacity reallocation.
% Structural model quantifies welfare trade-offs and enables counterfactuals.
\end{frame}


% ══════════════════════════════════════════════════════════
% POLICY AND SETTING
% ══════════════════════════════════════════════════════════
\section{Setting}

% ── Reform Overview ──────────────────────────────────────
\begin{frame}{Compra Ágil Reform (Ley 21.634)}
\begin{columns}[T]
  \column{0.49\textwidth}
  \textbf{Before (until Dec 11, 2024)}
  \begin{itemize}
    \item Compra Ágil limit: \textbf{30 UTM} ($\approx$ CLP 2M, USD 2k)
    \item Projects 30--100 UTM went through competitive licitaciones
    \item No explicit SME/local priority at first stage
    \item Non-local firms competed freely in all markets
  \end{itemize}
  \column{0.49\textwidth}
  \textbf{After (from Dec 12, 2024)}
  \begin{itemize}
    \item Compra Ágil limit: \textbf{100 UTM} ($\approx$ CLP 7M, USD 7k)
    \item Projects 30--100 UTM now use simplified process
    \item \textbf{SME and local firm priority} in first instance
    \item Non-local firms effectively excluded from this band
  \end{itemize}
\end{columns}

\bigskip
\textbf{Timeline:} Law 21.634 published November 2024; effective December 12, 2024.
First post-reform month with full exposure: January 2025.

\bigskip
\textbf{Scale:} The 30--100 UTM band accounts for $\approx$\,25--35\% of all
municipal procurement tenders by volume. The reform substantially narrows the
set of markets where non-local firms can participate.
\end{frame}

% ── Reform Impact on Procurement Mode ────────────────────
\begin{frame}{Projects Shifted Almost Entirely to Compra Ágil}
\centering
\includegraphics[width=0.96\linewidth]{ca_D4_share_compra_agil_by_utm_bucket.png}

\vspace{0.5em}
\small\textbf{Takeaway:} Nearly all 30--100 UTM projects switched from
licitaciones to Compra Ágil after the reform --- compliance is near-complete.
\end{frame}


\begin{frame}{Bidder Composition by Band}
\begin{columns}[c]
  \column{0.49\textwidth}\centering
  \textbf{\small 30--100 UTM}\\[2pt]
  \includegraphics[width=\linewidth]{ca_D8_bidder_types_avg_count_30_100utm_month.png}
  \column{0.49\textwidth}\centering
  \textbf{\small 100--500 UTM}\\[2pt]
  \includegraphics[width=\linewidth]{ca_D10_bidder_types_avg_count_100_500utm_month.png}
\end{columns}
\end{frame}

% ── Data ─────────────────────────────────────────────────
\begin{frame}{Data}
\begin{columns}[T]
  \column{0.5\textwidth}
  \textbf{ChileCompra (Mercado Público)}
  \begin{itemize}
    \item Universe of public procurement contracts in Chile
    \item Tender-level: estimated value (UTM), procurement type,
          region, sector, date
    \item Bid-level: firm RUT, bid amount, winner indicator
    \item \textbf{Coverage:} Jan 2022 -- Mar 2025
    \item \textbf{Sample:} Municipalidades sector;
          treated band (30--100 UTM);
          control bands (0--30, 100--200 UTM)
  \end{itemize}

  \column{0.5\textwidth}
  \textbf{Servicio de Impuestos Internos (SII)}
  \begin{itemize}
    \item Chilean Internal Revenue Service administrative data
    \item Firm-level: size classification (micro, small, medium, large),
          sector, location, annual sales
    \item \textbf{Use:} Classify bidders as SME (MiPYME) vs.\ large;
          local vs.\ non-local (matched to buyer region)
    \item \textbf{Coverage:} Annual data, matched to ChileCompra via RUT
  \end{itemize}
\end{columns}

\bigskip
\textbf{Final sample (Municipalidades):} $\approx$\,2,500 procuring entities;
$\approx$\,180,000 tenders; $\approx$\,650,000 bids over the study period.
\end{frame}

% ── Summary Statistics: Procurement ──────────────────────
\begin{frame}{Sample Overview: Pre-Reform Descriptives}
\textbf{[PLACEHOLDER — fill with summary statistics table]}

\medskip
Suggested content:
\begin{itemize}
  \item Table: N tenders, N bidders (mean), SME win share, log price ratio
        --- split by sector (Municipalidades vs.\ Obras Públicas) and UTM band
        (0--30, 30--100, 100--200 UTM), pre-reform only.
  \item Highlight that the treated band (30--100 UTM) had systematically
        lower bidder counts and higher non-local shares than control bands.
  \item Source: \texttt{did\_tender\_sample.parquet} + \texttt{did\_bid\_sample.parquet}.
\end{itemize}
\end{frame}

% ── Summary Statistics: SII ──────────────────────────────
\begin{frame}{Sample Overview: Bidder Characteristics (SII Data)}
\textbf{[PLACEHOLDER — fill with SII summary statistics]}

\medskip
Suggested content:
\begin{itemize}
  \item Share of bids by SME vs.\ large firm, by UTM band and sector
  \item Share of bids by local vs.\ non-local firm (home region = buyer region)
  \item Average firm size (sales, employees) of winning vs.\ losing bidders
  \item Geographic concentration: fraction of bidders from RM bidding in
        other regions (``export share'')
\end{itemize}
\end{frame}


% ══════════════════════════════════════════════════════════
% EMPIRICS
% ══════════════════════════════════════════════════════════
\section{Empirical Strategy \& Results}

% ── Facts About Local/Non-local Firms ────────────────────
\begin{frame}{Descriptive Facts: Local vs.\ Non-Local Bidders}
\textbf{[PLACEHOLDER --- fill with descriptive facts about local and non-local firm bids]}

\medskip
Suggested content:
\begin{itemize}
  \item Bar/line chart: share of bids by local vs.\ non-local firms by
        UTM band and region, pre-reform
  \item Map or bar chart: which regions send firms to other regions most?
        (foreshadows the export share analysis)
  \item Heterogeneity: non-local share is higher in remote regions,
        lower in RM (the inverse-U pattern from the binscatter)
\end{itemize}
\end{frame}

% ── DiD Specification ────────────────────────────────────
\begin{frame}{Difference-in-Differences: Specification}
\textbf{Estimating equation (TWFE):}
\[
  y_{it} = \alpha_i + \gamma_t + \beta\,\mathbf{1}[\text{treated}_i \times \text{post}_t] + \varepsilon_{it}
\]

\medskip
\textbf{Treatment and control groups} (defined by estimated contract value in UTM):
\begin{columns}[T]
  \column{0.48\textwidth}
  \begin{itemize}
    \item \textbf{Treated:} 30--100 UTM \\
          \small Tenders that switch from competitive bidding (licitación)
          to Compra Ágil after Dec 12, 2024.
  \end{itemize}
  \column{0.52\textwidth}
  \begin{itemize}
    \item \textbf{Control 1} (low): 1--30 UTM \\
          \small Already under Compra Ágil before the reform.
    \item \textbf{Control 2} (high): 100--200 UTM \\
          \small Above the new threshold; procurement process unchanged.
  \end{itemize}
\end{columns}

\medskip
\textbf{Fixed effects:} Procuring entity $\alpha_i$ + year-month $\gamma_t$.\\
\textbf{Standard errors:} Clustered by procuring entity.\\
\textbf{Sample:} Jan 2023 -- Mar 2025; Municipalidades sector.\\[4pt]
\small Results show control group: \textit{Control 2 (100--200 UTM)}.
\end{frame}

% ── IV-DiD Specification ─────────────────────────────────
\begin{frame}{IV-DiD: Specification}
\textbf{Motivation:} Intent-to-treat (DiD) estimates may under-state effects if
compliance is imperfect. We instrument actual Compra Ágil use with reform
eligibility.

\bigskip
\textbf{First stage (within-entity, within-month):}
\[
  \widetilde{CA}_{it} = \pi\,\widetilde{Z}_{it} + \nu_{it},
  \qquad
  Z_{it} = \mathbf{1}[\text{treated}_i \times \text{post}_t]
\]

\textbf{Second stage:}
\[
  \tilde{y}_{it} = \beta_{\text{IV}}\,\widehat{\widetilde{CA}}_{it} + \varepsilon_{it}
\]

\begin{itemize}
  \item \textbf{Endogenous:} $CA_{it}$ --- indicator, tender processed via Compra Ágil
  \item \textbf{Instrument:} $Z_{it} = \text{Treated}_i \times \text{Post}_t$ --- reform eligibility
  \item Same FE and clustering as OLS DiD; Control 2 only (100--200 UTM)
  \item $\beta_{\text{IV}}$ is a LATE: effect of Compra Ágil \emph{use}
\end{itemize}
\end{frame}

% ── IV First Stage ───────────────────────────────────────
\begin{frame}{IV-DiD: First Stage (Municipalidades)}
\small
\medskip
\input{../output/did/tables/did_table_first_stage_munic.tex}
\end{frame}

% ── Event Studies ─────────────────────────────────────────
\begin{frame}{Event Studies: N Bidders and SME Share (Municipalidades)}
\begin{columns}[c]
  \column{0.5\textwidth}\centering
  \includegraphics[width=\linewidth]{../output/did/figures/diag_event_study_n_bidders_munic.png}
  \column{0.5\textwidth}\centering
  \includegraphics[width=\linewidth]{../output/did/figures/diag_event_study_sme_share_sii_munic.png}
\end{columns}
\end{frame}

\begin{frame}
      \centering 
      \includegraphics[width=.8\linewidth]{../output/did/figures/diag_event_study_winner_is_sme_sii_munic.png}
\end{frame}

% \begin{frame}{Event Studies: SME Winners and Win Price (Municipalidades)}
% \begin{columns}[c]
%   \column{0.5\textwidth}\centering
%   \includegraphics[width=\linewidth]{../output/did/figures/diag_event_study_winner_is_sme_sii_munic.png}
%   \column{0.5\textwidth}\centering
%   \includegraphics[width=\linewidth]{../output/did/figures/diag_event_study_log_win_price_ratio_munic.png}
% \end{columns}
% \end{frame}

% \begin{frame}{Event Study: Bid-Level Price (Municipalidades)}
% \centering
% \includegraphics[width=0.58\linewidth]{../output/did/figures/diag_event_study_log_sub_price_ratio_munic.png}

% \medskip
% \small\textit{Flat pre-trends; significant post-reform increase in winning and
% submitted prices --- consistent with reduced competition (fewer non-local firms)
% not fully offset by local entry.}
% \end{frame}

% ── IV-DiD Results ───────────────────────────────────────
\begin{frame}{IV-DiD: Entry at Bidding Stage (Panel A, Municipalidades)}
\small\medskip
\input{../output/did/tables/did_table_compare_iv_A_munic.tex}
\end{frame}

\begin{frame}{IV-DiD: Bidder Composition (Panel B, Municipalidades)}
\small\medskip
\input{../output/did/tables/did_table_compare_iv_B_munic.tex}
\end{frame}

\begin{frame}{IV-DiD: Winner Characteristics (Panel C, Municipalidades)}
\small\medskip
\input{../output/did/tables/did_table_compare_iv_C_munic.tex}
\end{frame}

\begin{frame}{IV-DiD: Bid Outcomes (Panel D, Municipalidades)}
\small\medskip
\input{../output/did/tables/did_table_compare_iv_D_munic.tex}
\end{frame}

% ── Validity Checks ───────────────────────────────────────
% \begin{frame}{DiD Validity Checks}
% \textbf{Three diagnostic tests (Control 2 = 100--200 UTM):}
% \begin{enumerate}
%   \item \textbf{Wald pre-trend F-test} --- joint test
%         $H_0: \beta_{-8} = \cdots = \beta_{-2} = 0$
%         using full cluster-robust covariance matrix of event-study coefficients
%   \item \textbf{Time placebo} --- DiD run on pre-reform data only
%         ($k \in [-8,-1]$), fake reform at $k = -4$; coefficient should be
%         $\approx 0$
%   \item \textbf{Pre-reform balance} --- OLS of entity-level pre-reform means
%         on Treated; tests level differences before reform
% \end{enumerate}
% \end{frame}

% \begin{frame}{DiD: Pre-Trend Wald Test (Municipalidades)}
% \small\medskip
% \input{../output/did/tables/diag_pretrend_wald_munic.tex}
% \end{frame}

% \begin{frame}{DiD: Time Placebo Test (Municipalidades)}
% \small\medskip
% \input{../output/did/tables/diag_placebo_munic.tex}
% \end{frame}

% \begin{frame}{DiD: Pre-Reform Balance (Municipalidades)}
% \small\medskip
% \input{../output/did/tables/diag_balance_munic.tex}
% \end{frame}


% ══════════════════════════════════════════════════════════
% REGIONAL HETEROGENEITY
% ══════════════════════════════════════════════════════════
\section{Regional Heterogeneity}

\begin{frame}{Regional Heterogeneity: Market Structure Varies by Region}
\begin{columns}[c]
  \column{0.52\textwidth}
  \centering
  \includegraphics[width=\linewidth]{binscatter_q_nonlocal_munic.png}
  \column{0.48\textwidth}
  \textbf{Four pre-reform moderators (Municipalidades, 30--100 UTM):}
  \begin{itemize}
    \item \textbf{Nonlocal share:} fraction of bids by out-of-region firms
    \item \textbf{q\_pre:} avg.\ monthly tender count
    \item \textbf{n\_pot\_local:} unique local suppliers
    \item \textbf{Distance from Santiago:} geographic isolation (km)
  \end{itemize}
  \medskip
  \textit{Thick markets (high q\_pre) have lower non-local penetration ---
  market size proxies for local competitive capacity.}
\end{columns}
\end{frame}

\begin{frame}{Regional Heterogeneity: Specification}
\textbf{Question:} Does the effect of Compra Ágil vary with pre-reform
regional characteristics?

\bigskip
\textbf{Estimating equation (OLS DiD):}
\[
  y_{irt} = \alpha_i + \gamma_t
    + \beta_1\,(T_i \times \text{Post}_t)
    + \beta_2\,(T_i \times \text{Post}_t \times Z_r)
    + \varepsilon_{irt}
\]

\medskip
\begin{itemize}
  \item $Z_r$ = standardised pre-reform regional moderator
        (one specification per moderator)
  \item $\beta_1$ = average treatment effect;
        \quad $\beta_2$ = heterogeneity w.r.t.\ $Z_r$
  \item Same FE and clustering as baseline DiD
        (entity + year-month FE; SEs by PE)
  \item \textbf{Sample:} Municipalidades sector
\end{itemize}

\medskip
\textbf{Five moderators $Z_r$:} non-local share, tender count, tender value,
potential local firms, distance from Santiago.
\end{frame}

\begin{frame}{Interacted DiD --- Non-local Bidder Share}
\scriptsize\medskip
\resizebox{\textwidth}{!}{%
\input{../output/did/tables/hetero_interacted_nonlocal_share_pre_munic.tex}}
\end{frame}

\begin{frame}{Interacted DiD --- Potential Local Firms}
\small\medskip
\resizebox{\textwidth}{!}{%
\input{../output/did/tables/hetero_interacted_n_pot_local_munic.tex}}
\end{frame}

\begin{frame}{Interacted DiD --- Distance from Santiago}
\small\medskip
\resizebox{\textwidth}{!}{%
\input{../output/did/tables/hetero_interacted_dist_from_santiago_munic.tex}}
\end{frame}

\begin{frame}{Interacted DiD --- Total Value of Procurement}
\small\medskip
\resizebox{\textwidth}{!}{%
\input{../output/did/tables/hetero_interacted_totval_pre_munic.tex}}
\end{frame}

% \begin{frame}{Heterogeneity: N Local and N Nonlocal Bidders}
% \begin{columns}[c]
%   \column{0.5\textwidth}\centering
%   \vspace{.2cm} \\ 
%   \textbf{\small N Local Bidders}\\
%       \vspace{.2cm}
%   \includegraphics[width=0.7\linewidth]{hetero_coefplot_n_local_nonlocal_share_pre_munic.png}
%   \column{0.5\textwidth}\centering
%   \vspace{.2cm} \\ 
%   \textbf{\small N Non-local Bidders}\\
%   \vspace{.2cm}
%   \includegraphics[width=0.7\linewidth]{hetero_coefplot_n_nonlocal_nonlocal_share_pre_munic.png}
% \end{columns}
% \medskip
% \small\textit{Regions with high pre-reform non-local share see the largest local
% entry gains.}
% \end{frame}

% \begin{frame}{Heterogeneity: Total Bidders and SME Share (by Distance)}
% \begin{columns}[c]
      
%   \column{0.5\textwidth}\centering
%   \vspace{.2cm} \\ 
%   \textbf{\small N Bidders}\\
%   \vspace{.2cm}
%   \includegraphics[width=0.7\linewidth]{hetero_coefplot_n_bidders_dist_from_santiago_munic.png}
%   \column{0.5\textwidth}\centering
%   \vspace{.2cm}\\
%   \textbf{\small SME Share (SII)}\\
%   \vspace{.2cm}
%   \includegraphics[width=0.7\linewidth]{hetero_coefplot_sme_share_sii_dist_from_santiago_munic.png}
% \end{columns}
% \medskip
% \small\textit{Distance from Santiago is an exogenous geographic instrument for
% pre-reform non-local penetration.}
% \end{frame}

% ══════════════════════════════════════════════════════════
% CROSS-MARKET SPILLOVERS
% ══════════════════════════════════════════════════════════
\section{Cross-Market Spillovers}

% \begin{frame}{Cross-Market Spillovers: Mechanism}
% \textbf{Research question:} Did firms freed from 30--100 UTM licitaciones
% redirect capacity to \textit{above-threshold} (100--500 UTM) tenders in their
% home regions, increasing competition there?

% \bigskip
% \textbf{Mechanism (spatial auction model):}
% \begin{enumerate}
%   \item Reform excludes non-local firms from 30--100 UTM Compra Ágil.
%   \item Firms from ``export regions'' (e.g.\ RM) had many cross-region bids
%         in that band before the reform --- these bids are now blocked.
%   \item Freed capacity $\Rightarrow$ these firms bid more in their home
%         region's above-threshold (100--500 UTM) tenders.
%   \item Prediction: N bidders (esp.\ large firms) $\uparrow$ post-reform
%         in high-export-share regions.
% \end{enumerate}

% \bigskip
% \textbf{Treatment:} \textit{Export share} =
% fraction of pre-reform bids by home-region-$r$ firms placed in \textit{other}
% regions (30--100 UTM band). Time-invariant; variation is cross-regional.
% \end{frame}

% \begin{frame}{Export Share by Region (Pre-Reform)}
% \centering
% \includegraphics[width=0.50\linewidth]{spillover_export_share_munic.png}

% \medskip
% \small\textit{Región Metropolitana has the highest export share: Santiago firms
% dominated peripheral-region procurement in the treated band.
% Far south and north regions have low export share.}
% \end{frame}

\begin{frame}
      \begin{center}
            \textbf{Under construction...}
      \end{center}
\end{frame}


% \begin{frame}{Cross-Market Spillovers: Specification}
% \textbf{Event-study TWFE} on above-threshold (100--500 UTM) region--month panel:
% \[
%   y_{rt} = \alpha_r + \gamma_t
%     + \sum_{k \neq -1} \delta_k \,\bigl(\mathbf{1}[t{=}k] \times
%       \mathrm{ExportShare}_r\bigr) + \varepsilon_{rt}
% \]

% \begin{itemize}
%   \item $r$ = region (16 Chilean regions), $t$ = year-month.
%   \item $k$ = months relative to reform ($k=0$ = December 2024).
%   \item $\alpha_r$ = region FE; $\gamma_t$ = year-month FE.
%   \item Estimated for \textbf{continuous} (standardized) and
%         \textbf{binary} (top-tercile) export share.
%   \item Outcomes: N bidders, N local, N non-local, N large, N SME, Pr(single).
%   \item SE clustered at region level (16 clusters).
% \end{itemize}
% \end{frame}

% \begin{frame}{Spillovers: N Bidders and Pr(Single Bidder)}
% \begin{columns}[c]
%   \column{0.5\textwidth}\centering
%   \includegraphics[width=\linewidth]{spillover_event_study_n_bidders_munic.png}
%   \column{0.5\textwidth}\centering
%   \includegraphics[width=\linewidth]{spillover_event_study_single_bidder_munic.png}
% \end{columns}
% \medskip
% \small\textit{Total competition rises in high-export regions post-reform;
% single-bidder probability falls --- consistent with freed-capacity reallocation.}
% \end{frame}

% \begin{frame}{Spillovers: Local vs.\ Non-Local Bidders}
% \begin{columns}[c]
%   \column{0.5\textwidth}\centering
%   \includegraphics[width=\linewidth]{spillover_event_study_n_local_munic.png}
%   \column{0.5\textwidth}\centering
%   \includegraphics[width=\linewidth]{spillover_event_study_n_nonlocal_munic.png}
% \end{columns}
% \medskip
% \small\textit{Local bidder count rises in high-export regions (home firms
% re-enter above-threshold band). Non-local count adjusts as
% previously-exporting firms redirect to home markets.}
% \end{frame}

% \begin{frame}{Spillovers: Large Firms and Non-Local Share}
% \begin{columns}[c]
%   \column{0.5\textwidth}\centering
%   \includegraphics[width=\linewidth]{spillover_event_study_n_large_munic.png}
%   \column{0.5\textwidth}\centering
%   \includegraphics[width=\linewidth]{spillover_event_study_share_nonlocal_munic.png}
% \end{columns}
% \medskip
% \small\textit{Large-firm entry is the main driver of the spillover effect.
% Non-local share of above-threshold tenders rises in export regions as
% home-based large firms redirect bids.}
% \end{frame}

% \begin{frame}{Spillovers: SME Bidders}
% \centering
% \includegraphics[width=0.60\linewidth]{spillover_event_study_n_sme_munic.png}

% \medskip
% \small\textit{SME spillover effect is smaller --- SMEs tend to be local and
% were not the primary ``exporting'' firms. Effect concentrated among large firms.}
% \end{frame}


% ══════════════════════════════════════════════════════════
% STRUCTURAL MODEL
% ══════════════════════════════════════════════════════════
\section{Structural Model}

\begin{frame}
      \begin{center}
            \textbf{Under construction...}
      \end{center}
\end{frame}

% % ── Environment ──────────────────────────────────────────
% \begin{frame}{Model: Environment and Timing}
% \begin{columns}[T]
%   \column{0.5\textwidth}
%   \textbf{Regions and firms}
%   \begin{itemize}
%     \item $R$ regions with pairwise distances $d_{rs}$
%     \item Each period, region $r$ posts $Q_r$ procurement projects
%           at reserve price $\bar{p}$
%     \item Firms $i \in \mathcal{F}$ have home region $h_i$ and
%           capacity $K_i$ (max markets per period)
%     \item $N_r^{pot}$: potential local suppliers in region $r$
%   \end{itemize}

%   \column{0.5\textwidth}
%   \textbf{Timing}
%   \begin{enumerate}
%     \item \textbf{Demand} $\mathbf{Q}$ realized and observed
%     \item \textbf{Entry:} firm $i$ chooses set $M_i$,
%           $|M_i| \leq K_i$; pays entry cost $\kappa_r^i$
%     \item \textbf{Cost realization:} private cost drawn
%           for each entered market
%     \item \textbf{Bidding:} sealed-bid first-price auction;
%           lowest bidder wins
%   \end{enumerate}
% \end{columns}
% \end{frame}

% % ── Cost Structure ───────────────────────────────────────
% \begin{frame}{Model: Cost Structure}
% \textbf{Firm $i$'s cost for a project in region $r$:}
% \[
%   c_{ir} = \underbrace{\bar{c}}_{\text{baseline}}
%          + \underbrace{\delta \cdot d_{h_i, r}}_{\text{distance cost}}
%          - \underbrace{\lambda \cdot \mathbf{1}[h_i = r]}_{\text{local advantage}}
%          + \underbrace{\epsilon_{ir}}_{\text{private shock}}
% \]

% \medskip
% \begin{itemize}
%   \item $\bar{c} > 0$: common baseline project cost
%   \item $\delta > 0$: per-unit-distance cost (transport, logistics,
%         remote supervision)
%   \item $\lambda > 0$: local informational and relational advantage
%         (knowledge of local regulations, subcontractor networks,
%         permitting relationships)
%   \item $\epsilon_{ir} \sim F(\cdot)$ i.i.d.: idiosyncratic private cost shock
% \end{itemize}

% \medskip
% \textbf{Key:} A non-local firm from region $s \neq r$ faces an expected cost
% disadvantage of $\delta \cdot d_{sr} + \lambda$ relative to a local firm.
% Distance cost and informational advantage are \emph{separately} identified.
% \end{frame}

% % ── Entry Costs ──────────────────────────────────────────
% \begin{frame}{Model: Entry Costs and Capacity}
% \textbf{Entry cost for firm $i$ into market $r$:}
% \[
%   \kappa_r^i = \kappa_0 + \kappa_d \cdot d_{h_i, r}
% \]

% \begin{itemize}
%   \item $\kappa_0 > 0$: baseline participation cost (preparing bids,
%         learning about projects) --- same for all firms
%   \item $\kappa_d > 0$: additional cost of distant operations
%   \item Local firms pay only $\kappa_0$
% \end{itemize}

% \bigskip
% \textbf{Capacity constraint ($K_i$):}

% Firm $i$ can enter at most $K_i$ markets per period. This is the key ingredient
% that \textit{links markets}: a firm bidding in peripheral regions uses up
% capacity that could be deployed at home.

% \medskip
% Entry set choice:
% \[
%   \max_{M_i \subseteq \{1,\dots,R\},\; |M_i| \leq K_i}
%     \sum_{r \in M_i} \pi_{ir}(n_r^e)
% \]
% \end{frame}

% % ── Policy Environment ───────────────────────────────────
% \begin{frame}{Model: Policy Environment}
% \textbf{Pre-reform:} All firms may enter any market.

% \bigskip
% \textbf{Post-reform --- three tiers by project value:}

% \medskip
% \begin{enumerate}
%   \item \textbf{Strong preference (Compra Ágil, $< \underline{v}$ = 100 UTM):}\\
%         Only local firms + EMT may enter. Non-local firms \textbf{fully excluded}.

%   \vspace{0.5em}
%   \item \textbf{Weak preference ($\underline{v}$--$\bar{v}$, 100--500 UTM):}\\
%         All firms may enter, but local firms receive bid preference $\alpha > 0$.\\
%         Non-local firm $i$ wins only if $b_i < b_j - \alpha$ for all locals $j$.

%   \vspace{0.5em}
%   \item \textbf{No preference ($> \bar{v}$ = 500 UTM):}\\
%         Status quo; all firms compete on equal terms.
% \end{enumerate}

% \bigskip
% \small\textbf{Note:} Theoretical analysis focuses on the strong-form case
% (full exclusion, 30--100 UTM band) --- this is the empirically relevant tier
% and generates the starkest predictions.
% \end{frame}

% % ── Equilibrium ─────────────────────────────────────────
% \begin{frame}{Model: Equilibrium}
% \textbf{Bidding stage (conditional on $n_r$ entrants):}
% \[
%   E[p_r \mid n]
%     = \underbrace{E[c_{(1:n)}]}_{\text{efficient cost}}
%     + \underbrace{\frac{1}{n} E\!\left[\frac{F(c_{(1:n)})}{f(c_{(1:n)})}\right]}_{\text{markup (decreasing in }n)}
% \]
% More entrants $\Rightarrow$ lower prices via both lower $c_{(1:n)}$ and lower markup.

% \bigskip
% \textbf{Entry equilibrium:} vector $(n_1^*, \dots, n_R^*)$ such that
% for each market $r$ and each entering firm $i$:
% \[
%   Q_r \cdot \Pi_{ir}^{bid}(n_r^*) \geq \kappa_r^i
%   \quad \text{(with equality for marginal entrant)}
% \]

% \medskip
% \textbf{Multi-market:} the capacity constraint means
% \[
%   n_r^*\!\left(\{n_s^*\}_{s \neq r}\right) \quad \forall\; r
% \]
% must be solved as a system --- markets are interdependent through firm capacity.
% \end{frame}

% % ── Policy Effects: Direct ───────────────────────────────
% \begin{frame}{Model: Direct Effect of the Policy}
% \textbf{Immediate effect:} bidding pool in market $r$ shrinks from
% $n_r^{pre}$ to $n_r^L$ (local entrants only).

% \bigskip
% Government pays more if excluded non-local firms were cost-competitive.
% How much competition recovers depends on three factors:

% \medskip
% \begin{enumerate}
%   \item \textbf{$N_r^{pot}$ (local firm pool):} if large relative to equilibrium
%         entry, new entry can fully replace excluded firms. If small (remote regions),
%         it cannot.
%   \vspace{0.3em}
%   \item \textbf{$Q_r$ (demand volume):} thick demand sustains more entrants.
%         Thin, lumpy demand cannot support competitive local markets.
%   \vspace{0.3em}
%   \item \textbf{$\delta \cdot d_{sr} + \lambda$ (cost gap):} if excluded firms
%         were from nearby regions, the government loses genuinely competitive bidders.
%         If from far away, they were high-cost anyway.
% \end{enumerate}
% \end{frame}

% % ── Cross-Market Spillover ────────────────────────────────
% \begin{frame}{Model: Cross-Market Spillovers via Capacity}
% \textbf{Mechanism:} Consider Santiago-based firm $i$ that pre-reform entered
% both Santiago and rural region $r$ (using 2 of $K$ capacity units). Post-reform:

% \begin{enumerate}
%   \item Firm $i$ is excluded from $r$ --- recovers one capacity unit.
%   \item If there is another market $s$ with $\pi_{is}(n_s^e) > 0$ and $i$
%         was previously capacity-constrained, firm $i$ now enters $s$.
%   \item Otherwise, firm $i$ redirects capacity to Santiago's
%         above-threshold tenders.
% \end{enumerate}

% \bigskip
% \textbf{Effect on Santiago (export regions):}\\
% Redirected capacity increases entrants in above-threshold markets,
% \emph{reducing} procurement costs there.

% \medskip
% \textbf{Effect on rural markets (import regions):}\\
% Fewer non-local bidders in the 30--100 UTM band; competition falls,
% prices rise. But local entry partially offsets this.

% \bigskip
% \small\textit{This cross-market linkage is the key contribution of the
% multi-market framework --- single-market analyses miss it entirely.}
% \end{frame}

% ── Welfare ──────────────────────────────────────────────
% \begin{frame}{Model: Welfare}
% \textbf{Government procurement surplus:}
% \[
%   W_r^{gov} = Q_r \cdot \left(\bar{p} - E[p_r]\right)
% \]

% \textbf{Local economic benefit:}
% \[
%   W_r^{local} = Q_r \cdot \Pr[\text{local wins in } r] \cdot w_r
% \]
% where $w_r$ captures local employment, wages, and multiplier effects.

% \textbf{Total welfare:}
% \[
%   W = \sum_{r} \left[W_r^{gov} + \phi \cdot W_r^{local}\right]
% \]
% $\phi \geq 0$: social weight on local economic activity.

% \textbf{Change from policy:}
% \[
%   \Delta W = \sum_r Q_r
%     \left[
%       \underbrace{-(E[\tilde{p}_r] - E[p_r])}_{\text{cost effect}}
%       + \underbrace{\phi \cdot \Delta\!\Pr[\text{local wins}]_r \cdot w_r}_{\text{local benefit}}
%     \right]
% \]
% \textbf{Sign is ambiguous and varies across regions} --- the core tension.
% \end{frame}

% % ── Implications vs. Empirics ─────────────────────────────
% \begin{frame}{Model Implications vs.\ Current Empirical Results}
% \begin{columns}[T]
%   \column{0.5\textwidth}
%   \textbf{Matched by the data:}
%   \begin{itemize}
%     \item[\checkmark] SME bidder share $\uparrow$,
%           non-local share $\downarrow$ post-reform
%     \item[\checkmark] SME winners $\uparrow$,
%           large-firm winners $\downarrow$
%     \item[\checkmark] Win prices $\uparrow$ (cost effect $>$ competition recovery)
%     \item[\checkmark] Heterogeneity: regions with high pre-reform non-local
%           share see largest local entry gains
%     \item[\checkmark] Cross-market spillovers: above-threshold N bidders
%           $\uparrow$ in high-export regions (capacity reallocation confirmed)
%     \item[\checkmark] Large-firm spillover $>$ SME spillover
%   \end{itemize}

%   \column{0.5\textwidth}
%   \textbf{Unresolved / needs more work:}
%   \begin{itemize}
%     \item[\textbf{?}] Entry response timing: how quickly do new local
%           firms appear?
%     \item[\textbf{?}] $\lambda$ vs.\ $\delta$: cannot yet separate
%           distance cost from local informational advantage without
%           structural estimation
%     \item[\textbf{?}] $N_r^{pot}$ margin: limited data on potential
%           (non-entering) local firms
%     \item[\textbf{?}] Welfare weights $\phi$: distributional gains
%           require employment data (SII match)
%     \item[\textbf{?}] Bunching at thresholds: not yet tested
%   \end{itemize}
% \end{columns}
% \end{frame}


% ══════════════════════════════════════════════════════════
% ESTIMATION PLAN
% ══════════════════════════════════════════════════════════
% \section{Estimation Plan}

% \begin{frame}{Estimation Plan: Overview}
% \textbf{Goal:} Identify $(\delta, \lambda, \kappa_0, \kappa_d, K)$ using
% variation in entry and bidding across regions, distances, and the reform.

% \bigskip
% \textbf{Three-step strategy:}
% \begin{enumerate}
%   \item \textbf{Reduced-form moments} (already estimated):\\
%         Event-study and DiD coefficients provide model-consistent
%         moments relating entry changes to pre-reform market structure.

%   \vspace{0.3em}
%   \item \textbf{Bidding stage} (first, conditional on entry):\\
%         Estimate cost distribution $F(\cdot)$ from bid data using
%         standard first-price auction methods (e.g., GPV 2000;
%         Guerre, Perrigne, Vuong). Recovers $\delta$ and $\lambda$
%         from the difference in cost distributions for local vs.\
%         non-local firms at varying distances.

%   \vspace{0.3em}
%   \item \textbf{Entry stage} (second, using cost estimates):\\
%         Back out entry costs $(\kappa_0, \kappa_d)$ and capacity $K$
%         from the free-entry condition, using variation in $Q_r$,
%         $N_r^{pot}$, and the reform as instruments.
% \end{enumerate}
% \end{frame}

% \begin{frame}{Estimation Plan: Identifying Variation}
% \begin{columns}[T]
%   \column{0.5\textwidth}
%   \textbf{Identifying $\delta$ (distance cost):}
%   \begin{itemize}
%     \item Within-tender variation: same project, different bidder origins
%     \item Bid level $\sim$ distance of bidder's home region
%     \item Pre-reform data; no confound from exclusion
%     \item Controls: firm FE, tender FE, sector-region trends
%   \end{itemize}

%   \vspace{0.5em}
%   \textbf{Identifying $\lambda$ (local advantage):}
%   \begin{itemize}
%     \item Discontinuity: otherwise identical firms, local vs.\ non-local
%     \item Residual cost gap after controlling for $\delta \cdot d$
%     \item Validation: should predict pre-reform win share differential
%   \end{itemize}

%   \column{0.5\textwidth}
%   \textbf{Identifying $\kappa_0, \kappa_d$ (entry costs):}
%   \begin{itemize}
%     \item Cross-regional variation in entry rates vs.\ market size $Q_r$
%     \item Free-entry condition pins down $\kappa_0$ given cost estimates
%     \item $\kappa_d$: how entry by distant firms varies with distance
%           (pre-reform)
%   \end{itemize}

%   \vspace{0.5em}
%   \textbf{Identifying $K$ (capacity):}
%   \begin{itemize}
%     \item Post-reform spillovers: which firms enter more above-threshold
%           tenders after exclusion from 30--100 UTM?
%     \item Variation in export share $\times$ firm size
%     \item Direct measure: change in number of markets per firm
%   \end{itemize}
% \end{columns}
% \end{frame}

% ══════════════════════════════════════════════════════════
% COUNTERFACTUALS
% ══════════════════════════════════════════════════════════
% \section{Counterfactuals}

% \begin{frame}{Counterfactuals}
% \textbf{[PLACEHOLDER --- describe planned counterfactuals]}

% \medskip
% Once structural parameters are estimated, three main counterfactuals:

% \begin{enumerate}
%   \item \textbf{Optimal local preference threshold by region}\\
%         For each region $r$, find the threshold $\underline{v}_r^*$ that maximizes
%         total welfare $W_r^{gov} + \phi \cdot W_r^{local}$.
%         Quantify the cost of the uniform 100 UTM policy.

%   \vspace{0.4em}
%   \item \textbf{Welfare decomposition}\\
%         Separate the cost effect (government pays more) from the
%         local benefit effect (local firms win more) across regions.
%         Map the spatial distribution of winners and losers from the policy.

%   \vspace{0.4em}
%   \item \textbf{Alternative policy designs}\\
%         Bid preference $\alpha$ (weak form) vs.\ full exclusion (strong form).
%         Which achieves more local wins at lower cost to the government?
%         Optimal $\alpha$ by market thickness.
% \end{enumerate}
% \end{frame}


% ══════════════════════════════════════════════════════════
% NEXT STEPS
% ══════════════════════════════════════════════════════════
% \section{Next Steps}

% \begin{frame}{Next Steps}
% \begin{columns}[T]
%   \column{0.5\textwidth}
%   \textbf{Empirics (short run)}
%   \begin{itemize}
%     \item Add descriptive facts about local/non-local firm bids
%           (pre-reform patterns; geography of export activity)
%     \item Fill summary statistics tables (procurement + SII)
%     \item Extend event study window as more post-reform months arrive
%     \item Test for bunching at 100 UTM threshold post-reform
%     \item Track new local entrants (first-time bidders)
%           by region post-reform
%   \end{itemize}

%   \column{0.5\textwidth}
%   \textbf{Structural estimation}
%   \begin{itemize}
%     \item Match SII firm data to ChileCompra bidders (RUT linkage)
%     \item Estimate bidding-stage cost model (GPV approach);
%           recover $\delta, \lambda$ from local/non-local bid distributions
%     \item Estimate entry model; recover $\kappa_0, \kappa_d$
%           using pre-reform cross-regional variation
%     \item Use post-reform spillovers to identify capacity $K$
%     \item Validate: simulate model-predicted entry response and
%           compare to event-study estimates
%   \end{itemize}
% \end{columns}

% \bigskip
% \begin{itemize}
%   \item \textbf{Longer run:} obtain SII employment data for welfare
%         decomposition; compute local multipliers $w_r$ by region
%   \item \textbf{Writing:} working paper draft (target: Q3 2026)
% \end{itemize}
% \end{frame}


% ══════════════════════════════════════════════════════════
% APPENDIX
% ══════════════════════════════════════════════════════════

\begin{frame}
\centering
\textbf{APPENDIX}
\end{frame}

\appendix

% ── Monthly project counts and values ────────────────────

\begin{frame}{App: Monthly Project Counts by Band}
\begin{columns}[c]
  \column{0.49\textwidth}\centering
  \textbf{\small 0--30 UTM}\\[2pt]
  \includegraphics[width=\linewidth]{ca_D13_monthly_tender_counts_stacked_0_30utm.png}
  \column{0.49\textwidth}\centering
  \textbf{\small 30--100 UTM (Treated Band)}\\[2pt]
  \includegraphics[width=\linewidth]{ca_D13_monthly_tender_counts_stacked_30_100utm.png}
\end{columns}
\end{frame}

\begin{frame}{App: Monthly Project Counts --- Higher Bands}
\begin{columns}[c]
  \column{0.49\textwidth}\centering
  \textbf{\small 100--500 UTM}\\[2pt]
  \includegraphics[width=\linewidth]{ca_D13_monthly_tender_counts_stacked_100_500utm.png}
  \column{0.49\textwidth}\centering
  \textbf{\small 500+ UTM}\\[2pt]
  \includegraphics[width=\linewidth]{ca_D13_monthly_tender_counts_stacked_500_plus_utm.png}
\end{columns}
\end{frame}

\begin{frame}{App: Monthly Project Value by Band}
\begin{columns}[c]
  \column{0.49\textwidth}\centering
  \textbf{\small 0--30 UTM}\\[2pt]
  \includegraphics[width=\linewidth]{ca_D14_monthly_tender_value_stacked_0_30utm.png}
  \column{0.49\textwidth}\centering
  \textbf{\small 30--100 UTM (Treated Band)}\\[2pt]
  \includegraphics[width=\linewidth]{ca_D14_monthly_tender_value_stacked_30_100utm.png}
\end{columns}
\end{frame}

\begin{frame}{App: Monthly Project Value --- Higher Bands}
\begin{columns}[c]
  \column{0.49\textwidth}\centering
  \textbf{\small 100--500 UTM}\\[2pt]
  \includegraphics[width=\linewidth]{ca_D14_monthly_tender_value_stacked_100_500utm.png}
  \column{0.49\textwidth}\centering
  \textbf{\small 500+ UTM}\\[2pt]
  \includegraphics[width=\linewidth]{ca_D14_monthly_tender_value_stacked_500_plus_utm.png}
\end{columns}
\end{frame}

% ── Bidder types ─────────────────────────────────────────


\begin{frame}{App: Bidder Share by Band (30--100 UTM)}
\centering
\includegraphics[width=0.90\linewidth]{ca_D20_bidder_type_share_30_100utm_month.png}
\end{frame}

\begin{frame}{App: Winner Share by Band (30--100 UTM)}
\centering
\includegraphics[width=0.90\linewidth]{ca_D24_winner_type_share_30_100utm_month.png}
\end{frame}

% ── OLS DiD Results ──────────────────────────────────────

\begin{frame}{App: OLS DiD --- Entry (Panel A, Municipalidades)}
\small\medskip
\input{../output/did/tables/did_table_compare_A_munic.tex}
\end{frame}

\begin{frame}{App: OLS DiD --- Bidder Composition (Panel B, Municipalidades)}
\small\medskip
\input{../output/did/tables/did_table_compare_B_munic.tex}
\end{frame}

\begin{frame}{App: OLS DiD --- Winner Characteristics (Panel C, Municipalidades)}
\small\medskip
\input{../output/did/tables/did_table_compare_C_munic.tex}
\end{frame}

\begin{frame}{App: OLS DiD --- Bid Outcomes (Panel D, Municipalidades)}
\small\medskip
\input{../output/did/tables/did_table_compare_D_munic.tex}
\end{frame}

% ── Heterogeneity details ─────────────────────────────────

% \begin{frame}{App: Binscatter --- Market Thickness vs.\ Non-local Share}
% \begin{center}
% \includegraphics[width=0.90\linewidth]{binscatter_q_nonlocal_munic.png}
% \end{center}
% \small\textit{Thick markets (high q\_pre) have systematically fewer non-local
% bidders --- log scale removes Santiago leverage.}
% \end{frame}

% \begin{frame}{App: Binscatter --- Moderator Pair Matrix}
% \begin{center}
% \includegraphics[width=0.85\linewidth]{binscatter_moderator_pairs_munic.png}
% \end{center}
% \small\textit{q\_pre and totval\_pre are highly correlated. n\_pot\_local is
% a complementary channel. Santiago ($\star$) has extreme leverage on raw scale.}
% \end{frame}



% ── Option B robustness ───────────────────────────────────

% \begin{frame}{App: Robustness --- Dec 2024 Excluded (Option B)}
% \textbf{Motivation:} The reform took effect Dec 12. The $k=0$ month captures
% only partial exposure ($\approx 19/31$ days treated).

% \bigskip
% \textbf{Option B:} Drop $k=0$ entirely. Post-period begins at $k=1$ (Jan 2025).
% \end{frame}

% \begin{frame}{App: Event Studies --- Option B}
% \begin{columns}[c]
%   \column{0.5\textwidth}\centering
%   \includegraphics[width=\linewidth]{../output/did/figures/diag_event_study_n_bidders_optB.png}
%   \column{0.5\textwidth}\centering
%   \includegraphics[width=\linewidth]{../output/did/figures/diag_event_study_sme_share_sii_optB.png}
% \end{columns}
% \end{frame}

% \begin{frame}{App: Pre-Trend Wald Test (Option B)}
% \small\medskip
% \input{../output/did/tables/diag_pretrend_wald_optB.tex}
% \end{frame}

% ── Model comparative statics ─────────────────────────────

% \begin{frame}{App: When Does the Policy Reduce Procurement Costs?}
% \begin{columns}[T]
%   \column{0.5\textwidth}
%   \textbf{Policy more likely to \emph{reduce} costs when:}
%   \begin{enumerate}
%     \item High $N_r^{pot}$ relative to $Q_r$: large local firm pool
%           replaces excluded non-locals
%     \item High $\lambda$: non-locals had local disadvantage anyway;
%           their exclusion doesn't remove the best bidders
%     \item High $\delta \cdot d_{sr}$: non-locals came from far away
%           and were high-cost
%     \item Low $K$ (tight capacity): freed capacity redirects home,
%           generating competition spillovers
%   \end{enumerate}

%   \column{0.5\textwidth}
%   \textbf{Policy more likely to \emph{increase} costs when:}
%   \begin{enumerate}
%     \item Low $N_r^{pot}$ relative to $Q_r$: few local firms,
%           cannot replace competition
%     \item Low $\lambda$, low $d_{sr}$: non-locals were genuinely
%           competitive nearby firms
%     \item High $K$ (ample capacity): no meaningful cross-market
%           reallocation
%     \item Lumpy demand ($Q_r$ volatile): local market too thin to
%           sustain competition
%   \end{enumerate}
% \end{columns}
% \end{frame}

% ── Bunching ─────────────────────────────────────────────

\begin{frame}{App: Bunching Tests --- Motivation}
\textbf{What are we testing?} If agencies strategically size contracts to stay
below a UTM threshold, we should observe an excess density of contracts just
below that threshold.

\medskip
\textbf{Method:} Fit a 4th-order polynomial to the contract-value histogram,
excluding a $\pm 5$ UTM window around the threshold. Excess bunching ratio
$= (\text{actual} - \text{counterfactual})/\text{counterfactual}$ in the
window. A band test compares the below/above-threshold count ratio
pre- vs.\ post-reform using a $\chi^2$ test.

\medskip
\textbf{Two thresholds tested:}
\begin{itemize}
  \item \textbf{100 UTM} (main): the local preference threshold introduced
        by the Compra Ágil reform
  \item \textbf{30 UTM} (placebo): the prior Compra Ágil ceiling, now an
        interior point of the treated band
\end{itemize}
\end{frame}

\begin{frame}{App: Bunching --- Compra Ágil at 100 UTM}
\centering
\includegraphics[width=0.75\linewidth]{ca_bunching_100utm_pre_post_2025q1.png}

\smallskip
\small\textit{Excess ratio: $3.79\times$ pre-reform $\to$ $4.58\times$
post-reform.  Below/above ratio: $3.70 \to 9.43$
($\chi^2 = 289.1$, $p < 0.001$).}
\end{frame}

\begin{frame}{App: Bunching --- Licitaciones at 100 UTM (Placebo)}
\centering
\includegraphics[width=0.75\linewidth]{lic_bunching_100utm_pre_post_2025q1.png}

\smallskip
\small\textit{Excess ratio: $1.02\times$ pre-reform $\to$ $0.98\times$
post-reform---no strategic bunching.  Below/above ratio falls from
$5.05 \to 1.88$ ($\chi^2 = 710.5$, $p < 0.001$), reflecting Compra Ágil
absorbing sub-100 UTM volume rather than threshold manipulation.}
\end{frame}

\begin{frame}{App: Bunching at 30 UTM (Old Threshold)}
\begin{columns}[c]
  \column{0.49\textwidth}\centering
  \textbf{\small Compra Ágil}\\[4pt]
  \includegraphics[width=\linewidth]{ca_bunching_30utm_pre_post_2025q1.png}
  \column{0.49\textwidth}\centering
  \textbf{\small Licitaciones}\\[4pt]
  \includegraphics[width=\linewidth]{lic_bunching_30utm_pre_post_2025q1.png}
\end{columns}
\smallskip
\small\textit{Pre-reform: massive below-threshold bunching in Compra Ágil
(below/above $= 42.8$).  Post-reform: collapses to $2.79$
($\chi^2 = 74{,}299$, $p < 0.001$).  Bunching migrates from 30~UTM to
100~UTM after the reform.}
\end{frame}

\begin{frame}{App: Bunching --- Band Test Summary}
\small
\begin{center}
\begin{tabular}{llcccc}
\toprule
Threshold & Contract type & Pre below/above & Post below/above
          & $\chi^2$ & $p$-value \\
\midrule
100 UTM & Compra Ágil  & 3.70 & 9.43  & 289.1       & $<0.001$ \\
100 UTM & Licitaciones & 5.05 & 1.88  & 710.5       & $<0.001$ \\
30 UTM  & Compra Ágil  & 42.8 & 2.79  & $74{,}299$  & $<0.001$ \\
30 UTM  & Licitaciones & 0.19 & 1.69  & $1{,}530$   & $<0.001$ \\
\bottomrule
\end{tabular}
\end{center}
\medskip
\small\textit{All tests compare the below/above ratio pre- vs.\
post-2025Q1.  The Compra Ágil 30~UTM result is the sharpest signal:
extreme pre-reform bunching at the old ceiling disappears after the
threshold migrates to 100~UTM.}
\end{frame}

\end{document}
